% Documentation du projet A1 (filtrage collaboratif), groupe 40 « La piscine ».
% Partie 1 : les données, décrites à partir des fichiers. Partie 2 : les méthodes testées, dans l'ordre.
% Compilation (depuis docs/documentation/) : pdflatex documentation.tex (deux fois).
\documentclass[11pt,a4paper]{article}

\usepackage[T1]{fontenc}
\usepackage[utf8]{inputenc}
\usepackage[french]{babel}
\usepackage{lmodern}
\usepackage{microtype}
\usepackage{amsmath,amssymb}
\usepackage{booktabs}
\usepackage{tabularx}
\usepackage{xcolor}
\usepackage{enumitem}
\usepackage[margin=2.4cm]{geometry}
\usepackage[most]{tcolorbox}
\usepackage[hidelinks]{hyperref}

\definecolor{accent}{HTML}{2A6FDB}
\setlist{itemsep=2pt, topsep=4pt}
\setlength{\parskip}{5pt}
\setlength{\parindent}{0pt}
\setlength{\emergencystretch}{3em}
\renewcommand{\arraystretch}{1.2}

% Encadré « Définition » : un mot nouveau, défini avant d'être utilisé.
\newtcolorbox{definition}[1]{enhanced, breakable, colback=accent!5, colframe=accent,
  boxrule=0.6pt, arc=2pt, left=8pt, right=8pt, fonttitle=\bfseries, title={Définition : #1}}
% Bloc résultat d'une méthode.
\newcommand{\resultat}[1]{\par\textbf{Résultat sur le test :} #1\par}

\title{\textbf{Projet A1 --- Documentation}\\[4pt]
  \large Les données, puis les méthodes testées}
\author{Groupe 40 \og La piscine \fg{}}
\date{14 octobre 2026}

\begin{document}
\maketitle

Tout le contenu du dépôt (code, expériences, résultats, documents) a été généré par
Claude Code. Ce document décrit d'abord les données, à partir des fichiers eux-mêmes,
puis les méthodes testées, dans l'ordre où elles l'ont été. Chaque mot technique est
défini dans un encadré bleu avant d'être utilisé.

\tableofcontents

%=============================================================================
\part{Les données}
%=============================================================================

\section{Les trois fichiers}

Le cours fournit trois fichiers, dans le dossier \texttt{data/} :

\begin{center}
\begin{tabularx}{\linewidth}{@{}lX@{}}
\toprule
\textbf{Fichier} & \textbf{Contenu} \\
\midrule
\texttt{ratings\_train.npy} & un tableau de notes de films (les notes qu'on a le droit d'utiliser) \\
\texttt{ratings\_test.npy} & un second tableau, de même forme, avec d'autres notes (pour vérifier nos prédictions) \\
\texttt{namesngenre.npy} & le titre et les genres de chaque film \\
\bottomrule
\end{tabularx}
\end{center}

L'extension \texttt{.npy} désigne un fichier de la bibliothèque Python \emph{numpy} :
il contient un tableau de nombres (ou de textes), qu'on ouvre avec \texttt{np.load}.

\section{Le tableau de notes}

\subsection{Sa forme}

\texttt{ratings\_train.npy} est un tableau de \textbf{610 lignes} et \textbf{4\,980 colonnes}.

\begin{itemize}
  \item \textbf{Une ligne = un utilisateur} : une personne qui a noté des films. Les
        utilisateurs sont anonymes ; on les désigne par leur numéro de ligne, de 0 à 609.
  \item \textbf{Une colonne = un film}, numéroté de 0 à 4\,979. Son titre est donné par
        le fichier \texttt{namesngenre.npy} (section~\ref{sec:films}).
  \item \textbf{Une case} (ligne $u$, colonne $i$) contient la note que l'utilisateur $u$
        a donnée au film $i$, ou rien s'il ne l'a pas noté.
\end{itemize}

\begin{definition}{note}
Un nombre entre 0,5 et 5, par pas de 0,5 : 0,5 ; 1 ; 1,5 ; \dots ; 5. C'est le nombre
d'étoiles qu'une personne a données à un film.
\end{definition}

\begin{definition}{NaN}
\og Not a Number \fg. C'est la valeur que contient une case vide : elle signifie
\og cet utilisateur n'a pas noté ce film \fg (ou, dans le fichier d'entraînement, que la
note est cachée). Ce n'est pas une note de 0.
\end{definition}

\subsection{Un extrait réel}

Voici les cinq utilisateurs qui ont noté le plus de films, croisés avec quatre des films
les plus notés. \og --- \fg{} est une case vide (NaN).

\begin{center}
\small\setlength{\tabcolsep}{4pt}
\begin{tabular}{lcccc}
\toprule
 & Forrest Gump & Shawshank Redemption & Pulp Fiction & Silence of the Lambs \\
 & (film 256) & (film 227) & (film 211) & (film 413) \\
\midrule
Utilisateur 413 & 5 & --- & --- & --- \\
Utilisateur 473 & --- & --- & 4 & 4,5 \\
Utilisateur 598 & 3,5 & --- & 5 & 3 \\
Utilisateur 447 & --- & --- & --- & --- \\
Utilisateur 273 & 4,5 & --- & 5 & --- \\
\bottomrule
\end{tabular}
\end{center}

Même pour les utilisateurs les plus actifs et les films les plus connus, la plupart des
cases sont vides.

\subsection{Combien de cases sont remplies}

Le tableau a $610 \times 4\,980 = 3\,037\,800$ cases. Seules \textbf{31\,598} contiennent
une note, soit \textbf{1,04\,\%}. Les 99\,\% restants sont des NaN.

\begin{center}
\begin{tabular}{lrrr}
\toprule
 & \textbf{minimum} & \textbf{médiane} & \textbf{maximum} \\
\midrule
Notes par utilisateur & 2 & 24 & 815 \\
Notes par film & 1 & 3 & 109 \\
\bottomrule
\end{tabular}
\end{center}

La \emph{médiane} est la valeur du milieu : la moitié des films ont 3 notes ou moins.
Plus précisément, 57\,\% des films ont 3 notes ou moins, et 1\,301 films (26\,\%) n'en
ont qu'une seule. Chaque film a au moins une note.

\subsection{Les notes données}

La moyenne de toutes les notes est \textbf{3,52}. Les notes entières et hautes dominent :

\begin{center}
\begin{tabular}{lrrrrrrrrrr}
\toprule
Note & 0,5 & 1 & 1,5 & 2 & 2,5 & 3 & 3,5 & 4 & 4,5 & 5 \\
\midrule
Nombre & 382 & 853 & 526 & 2\,308 & 1\,665 & 6\,231 & 4\,153 & 8\,532 & 2\,675 & 4\,273 \\
\bottomrule
\end{tabular}
\end{center}

\section{Le tableau de test}

\texttt{ratings\_test.npy} a exactement la même forme (610 lignes, 4\,980 colonnes, les
mêmes utilisateurs et les mêmes films), et contient lui aussi 31\,598 notes. Mais
\textbf{aucune case n'est remplie dans les deux tableaux à la fois} : les notes du test
sont des notes cachées dans le tableau d'entraînement.

\textbf{Exemple réel.} L'utilisateur 0 a noté \emph{Grumpier Old Men (1995)} (film 2).
Dans \texttt{ratings\_train.npy}, la case (0, 2) est vide ; dans \texttt{ratings\_test.npy},
elle vaut 4. Une bonne méthode, qui ne voit que le tableau d'entraînement, doit prédire
une valeur proche de 4 pour cette case.

Une difficulté importante : 21\,\% des notes du test portent sur un film qui a moins de
5 notes dans le tableau d'entraînement. Pour ces films, on dispose de très peu
d'information.

\section{Le fichier des films}\label{sec:films}

\texttt{namesngenre.npy} est un tableau de 4\,980 lignes et 2 colonnes : la ligne $i$
décrit le film de la colonne $i$ des tableaux de notes. Ses cinq premières lignes :

\begin{center}
\begin{tabular}{rll}
\toprule
\textbf{Film} & \textbf{Titre (année)} & \textbf{Genres} \\
\midrule
0 & Toy Story (1995) & Adventure|Animation|Children|Comedy|Fantasy \\
1 & Jumanji (1995) & Adventure|Children|Fantasy \\
2 & Grumpier Old Men (1995) & Comedy|Romance \\
3 & Waiting to Exhale (1995) & Comedy|Drama|Romance \\
4 & Father of the Bride Part II (1995) & Comedy \\
\bottomrule
\end{tabular}
\end{center}

\begin{itemize}
  \item L'année est entre parenthèses à la fin du titre ; elle va de 1902 à 2018
        (absente pour 3 films).
  \item Les genres sont séparés par \og | \fg. Il y en a 20, de \emph{Drama}
        (2\,253 films) à \emph{Film-Noir} (50 films), plus \og (no genres listed) \fg
        pour 8 films.
\end{itemize}

\section{Ce qu'il faut produire, et comment c'est noté}

\subsection{La tâche}

Le script \texttt{generate.py} reçoit un tableau de notes troué et doit écrire, dans
\texttt{output.npy}, \textbf{le même tableau avec une valeur dans toutes les cases}.
La valeur mise dans une case vide est une \textbf{prédiction} : notre estimation de la
note que l'utilisateur donnerait à ce film.

Sur la plateforme du cours, \texttt{generate.py} reçoit l'entraînement et le test réunis
(63\,196 notes). Il est noté sur un troisième jeu de notes, caché, que nous n'avons pas.
Limites : 20 minutes, 4 c\oe urs de processeur, pas de carte graphique.

\subsection{La mesure de l'erreur}

\begin{definition}{RMSE (erreur quadratique moyenne)}
Pour chaque note cachée, on calcule l'écart entre la vraie note et la prédiction, on le
met au carré, on fait la moyenne de ces carrés, puis on prend la racine carrée :
\[
  \mathrm{RMSE} = \sqrt{\text{moyenne de } (\text{vraie note} - \text{prédiction})^2} .
\]
Le résultat s'exprime en étoiles. Le carré rend les écarts positifs et pénalise davantage
les grosses erreurs. \textbf{Plus la RMSE est basse, meilleure est la méthode.}
\end{definition}

Exemple : deux notes cachées, de vraies valeurs 4 et 3, prédites 3 et 3. Écarts : 1 et 0.
Carrés : 1 et 0. Moyenne : 0,5. Racine : 0,71. La RMSE vaut 0,71 étoile.

La plateforme affiche aussi l'\emph{accuracy}, le pourcentage de notes exactement
devinées une fois la prédiction arrondie à la demi-étoile, et le temps de calcul.

\subsection{Comment on évalue nos méthodes en local}

On n'a pas les notes cachées de la plateforme. On imite donc la situation avec nos deux
fichiers :

\begin{definition}{entraînement, validation, test}
\begin{itemize}[leftmargin=*]
  \item \textbf{Entraînement} : les notes que la méthode voit pour calculer ses prédictions.
  \item \textbf{Validation} : 10\,\% des notes d'entraînement, tirées au hasard (toujours
        les mêmes) et mises de côté. Elles servent à choisir les réglages d'une méthode :
        on essaie plusieurs réglages et on garde celui qui a la plus petite RMSE sur
        la validation.
  \item \textbf{Test} : le fichier \texttt{ratings\_test.npy}. Une fois le réglage choisi,
        la méthode est recalculée sur tout l'entraînement, puis sa RMSE est mesurée
        \textbf{une seule fois} sur le test. Le test ne sert jamais à choisir.
\end{itemize}
\end{definition}

\begin{definition}{réglage (hyperparamètre)}
Un nombre que l'on fixe soi-même avant de lancer une méthode, par exemple le nombre de
voisins à regarder. On le choisit en comparant plusieurs valeurs sur la validation.
\end{definition}

Tous les chiffres \og RMSE test \fg{} de ce document sont obtenus ainsi.

%=============================================================================
\part{Les méthodes testées, dans l'ordre}
%=============================================================================

Chaque méthode est un moyen de remplir les cases vides. Elles sont présentées dans
l'ordre où elles ont été testées. La plupart réutilisent la précédente.

\section{La moyenne globale}

\textbf{Ce qu'on calcule.} La moyenne des 31\,598 notes d'entraînement : 3,52. On met
3,52 dans toutes les cases vides.

\resultat{RMSE 1,037. C'est le point de départ : toute méthode utile doit faire mieux.}

\section{Les biais}\label{sec:biais}

\textbf{L'idée.} Certains utilisateurs notent plus haut que les autres ; certains films
sont mieux notés que les autres. On corrige la moyenne globale de ces deux effets.

\begin{definition}{biais d'un utilisateur, biais d'un film}
Le \textbf{biais d'un utilisateur} est le nombre qu'il faut ajouter à la moyenne globale
pour tenir compte de sa façon de noter : positif s'il note plus haut que les autres,
négatif s'il note plus bas. Le \textbf{biais d'un film} est la même chose pour un film :
positif s'il reçoit de meilleures notes que la moyenne, négatif sinon.
\end{definition}

La prédiction pour l'utilisateur $u$ et le film $i$ est une addition :
\[
  \text{prédiction} = \text{moyenne globale} + \text{biais de } u + \text{biais de } i .
\]

\textbf{Exemple réel : utilisateur 0, film 2 (\emph{Grumpier Old Men}).}
\begin{itemize}
  \item Moyenne globale : 3,52.
  \item L'utilisateur 0 a donné 71 notes, de moyenne 4,35, soit 0,83 au-dessus de 3,52.
        Son biais calculé est \textbf{+0,63} : un peu moins que 0,83, parce qu'il a
        surtout noté des films que tout le monde note bien ; une partie de ses notes
        hautes s'explique par les films, pas par lui.
  \item \emph{Grumpier Old Men} a 17 notes, de moyenne 3,03. Son biais est \textbf{$-0{,}36$}.
  \item Prédiction : $3{,}52 + 0{,}63 - 0{,}36 = \mathbf{3{,}80}$. La vraie note (dans le
        test) est 4.
\end{itemize}

\textbf{Le problème des films peu notés.} \emph{Party Girl (1995)} (film 133) n'a qu'une
note : un 5. Calculé naïvement, son biais serait d'environ $+1{,}5$, sur la foi d'une
seule personne. On ajoute donc une prudence :

\begin{definition}{régularisation}
Une correction qui rapproche de 0 les nombres calculés sur peu de données. Pour un biais :
on divise la somme des écarts par (nombre de notes $+ \lambda$) au lieu du nombre de notes.
$\lambda$ (\og lambda \fg) est un réglage. Avec $\lambda = 2$, un film à 1 note voit son
écart divisé par 3 ; un film à 100 notes, par 102 au lieu de 100, presque rien.
\end{definition}

Avec $\lambda = 2$, le biais de \emph{Party Girl} tombe à $+0{,}45$. Les biais sont
calculés en alternant : biais des films à biais des utilisateurs fixés, puis l'inverse,
vingt fois. Cela prend 0,01 seconde.

\resultat{RMSE \textbf{0,8746} (meilleur réglage : $\lambda = 2$). C'est la méthode de
référence.}

\begin{definition}{résidu}
L'écart entre une vraie note et la prédiction des biais. Pour l'utilisateur 0 et
\emph{Grumpier Old Men} : $4 - 3{,}80 = +0{,}20$. Le résidu, c'est ce que les biais
n'expliquent pas : le goût propre de cet utilisateur pour ce film. Toutes les méthodes
suivantes partent des biais et essaient de prédire les résidus.
\end{definition}

\section{La factorisation de matrice}\label{sec:mf}

C'est la méthode exigée par le sujet.

\textbf{L'idée.} On décrit chaque utilisateur par une courte liste de $k$ nombres, et
chaque film par une liste de $k$ nombres. Le résidu prédit pour un couple (utilisateur,
film) est obtenu en multipliant les deux listes terme à terme puis en additionnant.

\begin{definition}{facteurs, rang}
Les \textbf{facteurs} d'un utilisateur (ou d'un film) sont ses $k$ nombres. Personne ne
les choisit : la méthode les calcule pour que les prédictions collent aux notes connues.
Le \textbf{rang} est $k$, la longueur des listes (nous avons essayé de 1 à 50).
\end{definition}

Exemple avec $k = 2$ (valeurs inventées pour illustrer) : facteurs de l'utilisateur
$(0{,}8 \,;\, -0{,}5)$, facteurs du film $(0{,}9 \,;\, 0{,}1)$. Résidu prédit :
$0{,}8 \times 0{,}9 + (-0{,}5) \times 0{,}1 = 0{,}67$. Prédiction finale = moyenne + biais
de l'utilisateur + biais du film + 0,67.

Le nom vient de là : le tableau $610 \times 4\,980$ est approché par le produit d'un
tableau $610 \times k$ (les facteurs des utilisateurs) et d'un tableau $k \times 4\,980$
(ceux des films).

\textbf{Le danger.} Avec beaucoup de facteurs et peu de notes, la méthode peut reproduire
parfaitement les notes d'entraînement sans rien savoir prédire d'autre. On parle de
\emph{surapprentissage}. On ajoute donc une régularisation sur les facteurs (réglage
$\lambda_f$), comme pour les biais.

\textbf{Trois façons de calculer les facteurs, testées :}
\begin{itemize}
  \item \textbf{ALS} (moindres carrés alternés) : on fixe les facteurs des films et on
        calcule directement, par une formule, les meilleurs facteurs des utilisateurs ;
        puis l'inverse ; quinze fois.
  \item \textbf{ALS-WR} : la même chose, avec une régularisation plus forte pour les
        utilisateurs et les films qui ont beaucoup de notes.
  \item \textbf{SGD} (descente de gradient) : on parcourt les notes une par une et, à chaque
        fois, on modifie un peu les facteurs dans le sens qui réduit l'erreur sur cette note.
\end{itemize}

\resultat{ALS \textbf{0,8725} (rang 20, $\lambda_f = 10$) ; ALS-WR 0,8746 ; SGD 0,8775.
Gain faible sur les biais. Sans régularisation : erreur proche de 0 sur l'entraînement,
plus de 1,0 sur la validation. Avec une régularisation trop forte ($\lambda_f \geq 15$),
les facteurs valent 0 et on retrouve exactement les biais.}

\section{Soft-Impute}

\textbf{L'idée.} Même objectif que la factorisation, sans fixer $k$ à l'avance. On
travaille sur le tableau des résidus et on répète :
\begin{enumerate}
  \item remplir les cases vides avec la prédiction actuelle (0 au départ) ;
  \item \og lisser \fg{} le tableau complet, c'est-à-dire n'en garder que les grandes
        tendances communes à beaucoup d'utilisateurs et de films ;
  \item recommencer jusqu'à ce que le tableau ne change plus.
\end{enumerate}

\begin{definition}{SVD (décomposition en valeurs singulières)}
Une opération mathématique qui écrit n'importe quel tableau de nombres comme une somme de
\og motifs \fg{} simples, rangés du plus important au plus faible. L'importance d'un motif
est un nombre appelé \emph{valeur singulière}. Le lissage de Soft-Impute retire la même
quantité $\tau$ (\og tau \fg, un réglage) à chaque valeur singulière : les motifs faibles
tombent à 0 et disparaissent, les forts restent.
\end{definition}

\resultat{RMSE \textbf{0,8716} ($\tau = 10$), en 4,7 secondes. Avec $\tau \geq 20$, tous
les motifs disparaissent et on retrouve les biais.}

\section{Les plus proches voisins (kNN)}

\textbf{L'idée.} Pour prédire la note de l'utilisateur $u$ sur le film $i$, on regarde
comment $u$ a noté les films qui \og ressemblent \fg{} à $i$.

\begin{definition}{ressemblance entre deux films}
On prend les utilisateurs qui ont noté les deux films. Si ceux qui ont un résidu positif
sur l'un ont aussi un résidu positif sur l'autre (et pareil pour les négatifs), les films
se ressemblent. On le mesure par une \emph{corrélation}, un nombre entre $-1$ et $1$.
Comme une corrélation calculée sur 2 ou 3 personnes n'est pas fiable, on la réduit quand
peu d'utilisateurs ont noté les deux films (réglage \og shrink \fg).
\end{definition}

\textbf{Le calcul}, pour $u$ et $i$ :
\begin{enumerate}
  \item parmi les films notés par $u$, garder les $k$ qui ressemblent le plus à $i$
        (les \og voisins \fg ; réglage $k$) ;
  \item faire la moyenne des résidus de $u$ sur ces voisins, chacun pesant selon sa
        ressemblance avec $i$ ;
  \item ajouter cette moyenne à la prédiction des biais.
\end{enumerate}
Ajout de notre part : un petit nombre $\beta$ (\og bêta \fg) ajouté au dénominateur de la
moyenne. Il réduit la correction quand les voisins sont peu nombreux ou peu ressemblants.
Sur la validation : 0,908 sans $\beta$, 0,885 avec $\beta = 0{,}2$.

\resultat{RMSE \textbf{0,8705} ($k = 40$, shrink 100, $\beta = 0{,}2$). Meilleur sur les
films très notés ; légèrement moins bon que les biais sur les films à 1 ou 2 notes.}

\section{Les biais complétés par les genres}

\textbf{L'idée.} Utiliser le fichier des films. Avec les biais simples, la régularisation
rapproche de 0 le biais d'un film peu noté. Ici, on le rapproche plutôt de la valeur
attendue d'après ses genres, sa décennie de sortie et son nombre de notes. Cette valeur
attendue est apprise sur l'ensemble des films (une \emph{régression linéaire} : un poids
par genre et par décennie, qu'on additionne).

\resultat{RMSE \textbf{0,8736}. Gain global minime, mais c'est la seule méthode jusqu'ici
qui fait mieux que les biais sur les films à 1 ou 2 notes (0,940 contre 0,944).}

\section{Le mélange}

\textbf{L'idée.} Combiner les prédictions des cinq méthodes précédentes (biais,
biais + genres, Soft-Impute, kNN, ALS) par une somme pondérée :
\[
  \text{prédiction} = a_1 \times \text{prédiction}_1 + \dots + a_5 \times \text{prédiction}_5 + \text{constante}.
\]
Les poids $a_1, \dots, a_5$ sont choisis pour que la RMSE soit la plus petite possible sur
la validation (régression \emph{ridge} : une régression linéaire avec régularisation).
Poids obtenus : kNN 0,94 ; biais + genres 0,72 ; ALS 0,29 ; Soft-Impute 0,06 ;
biais $-1{,}01$. Le poids négatif des biais compense le fait qu'ils sont déjà contenus dans
les quatre autres méthodes.

\resultat{RMSE \textbf{0,8683}, mieux que chaque méthode seule.}

\section{AutoRec (réseau de neurones)}

\begin{definition}{réseau de neurones, auto-encodeur}
Un \textbf{réseau de neurones} est une fonction composée de couches de calculs simples
(multiplications, additions, puis une fonction qui \og écrase \fg{} les valeurs) ; ses
nombreux coefficients sont ajustés pour réduire l'erreur. Un \textbf{auto-encodeur} est un
réseau qui reçoit une liste de nombres, la résume en une liste plus courte, puis essaie de
reconstruire la liste de départ à partir de ce résumé.
\end{definition}

\textbf{Ici.} L'entrée est une colonne du tableau des résidus (les 610 résidus d'un film,
avec des 0 pour les cases vides). Le réseau la résume en 500 nombres, puis reconstruit
les 610 valeurs. Il est entraîné à bien reconstruire les cases connues ; les valeurs qu'il
reconstruit dans les cases vides sont les prédictions. Variantes testées : par film ou par
utilisateur, sur les notes ou sur les résidus, 100 à 500 nombres de résumé, 27 réglages.

\resultat{RMSE \textbf{0,8722}, pas mieux que les méthodes simples, et aucun gain dans le
mélange. Le réseau a plusieurs centaines de milliers de coefficients pour 31\,598 notes :
il surapprend.}

\section{Les méthodes bayésiennes}

\textbf{L'idée.} Toutes les méthodes précédentes calculent \emph{une} solution (un jeu de
biais et de facteurs). Avec si peu de notes, beaucoup de solutions différentes sont presque
aussi plausibles. Les méthodes bayésiennes en tirent des centaines au hasard, chacune
compatible avec les données, et font \textbf{la moyenne de leurs prédictions}.

\begin{definition}{échantillonnage de Gibbs}
La façon de tirer ces solutions : comme l'ALS, on alterne utilisateurs et films, mais au lieu
de prendre les meilleurs facteurs, on en tire au hasard selon leur plausibilité. La force
de la régularisation est elle-même tirée au hasard : il n'y a plus de $\lambda$ à régler.
Les premiers tirages, encore mauvais, sont jetés.
\end{definition}

Deux méthodes testées :
\begin{itemize}
  \item \textbf{BPMF} : la factorisation de matrice (biais + facteurs) en version
        bayésienne. 300 tirages moyennés.
  \item \textbf{FM bayésienne} (\emph{factorization machine}) : la même idée, en donnant
        en plus à la méthode, pour chaque note, la liste des films notés par l'utilisateur,
        la liste des utilisateurs qui ont noté le film, et les genres et la décennie du film.
        2\,000 tirages. Calculée avec la bibliothèque \texttt{myfm}.
\end{itemize}

Observations : un tirage seul a une RMSE d'environ 1,03 (pire que la moyenne globale) ;
c'est la moyenne des tirages qui est bonne. Parmi les informations ajoutées à la FM, la
plus utile est la liste des utilisateurs qui ont noté le film (validation : 0,8887 sans
rien, 0,8762 avec elle seule, 0,8728 avec les trois).

\resultat{BPMF \textbf{0,8646} (44 s) ; FM bayésienne \textbf{0,8552} (60 s ; environ
3 min sur entraînement + test). C'est le meilleur modèle seul, et le premier gain net sur
les biais ($-0{,}019$).}

\textbf{Ce qui n'a pas aidé la FM} (expériences \texttt{experiments/run\_fm\_variants.py}) :
la version \emph{ordinale} (prédire l'un des 10 niveaux de note au lieu d'un nombre :
validation 0,888 contre 0,874), donner à la FM les notes elles-mêmes et pas seulement
\og qui a noté quoi \fg{} (0,977 : elle apprend par c\oe ur), un rang plus grand, plus de
tirages, ou la moyenne de 3 chaînes de tirages indépendantes (test 0,8545 au lieu de
0,8552, pour 3 fois plus de temps).

\section{Autres méthodes de la littérature}

Testées pour vérifier qu'aucune famille de méthodes connue ne faisait nettement mieux
(\texttt{experiments/run\_literature\_models.py}, \texttt{experiments/run\_residual\_models.py}) :
\begin{itemize}
  \item \textbf{GC-MC} (réseau de neurones sur le graphe utilisateurs--films) : test 0,8792.
  \item \textbf{GLocal-K} (auto-encodeur à noyau, très bon sur MovieLens-100K) : test 0,8965.
  \item \textbf{EASE} (régression linéaire film--film, en une formule) sur les résidus des
        biais ou de la FM : aucun gain (validation 0,8896 et 0,8727).
  \item \textbf{kNN sur les erreurs de la FM} : aucun gain (0,8556 sur le test).
\end{itemize}
Sur si peu de notes, les réseaux de neurones font moins bien que la FM bayésienne.

\section{L'ensemble final}\label{sec:stack}

\textbf{L'idée.} Comme pour le mélange de la section précédente, on combine plusieurs
prédictions par une somme pondérée, mais cette fois avec la FM bayésienne. On a besoin
d'une nouvelle brique :

\begin{definition}{correction par arbres de décision (GBDT)}
Un petit modèle (une somme de 300 arbres de décision peu profonds) apprend l'\emph{erreur}
de la FM à partir de quelques indices : la prédiction de la FM, et pour l'utilisateur et le
film, leur nombre de notes, leur note moyenne et l'écart-type de leurs notes. Il apprend
par exemple que la FM se trompe d'une certaine façon sur les films à 1 ou 2 notes.
Pour apprendre ces erreurs honnêtement, la FM est recalculée 5 fois, chaque fois sans un
cinquième des notes, et ses erreurs sont mesurées sur ce cinquième (\emph{out-of-fold}).
\end{definition}

L'ensemble final combine quatre prédictions, avec des poids appris sur la validation
(régression ridge) :
\[
  \text{prédiction} = -0{,}25 + 0{,}53\,\text{FM corrigée} + 0{,}27\,\text{FM}_2
  + 0{,}25\,k\text{NN} + 0{,}01\,\text{FM}
\]
où FM$_2$ est une seconde FM qui ne connaît que \og qui a noté quoi \fg{} (sans les genres).
Le $k$NN est utile parce qu'il est bon là où la FM l'est moins : sur les films très notés.

Combinaisons comparées sur la validation (5 plis, \texttt{experiments/run\_stack.py}) :
FM seule 0,8726 ; FM + $k$NN + FM$_2$ 0,8656 ; FM corrigée 0,8662 ; FM corrigée + $k$NN
0,8645 ; les quatre 0,8636.

\resultat{RMSE \textbf{0,8470}, accuracy 0,253. Gain de $-0{,}008$ sur la FM seule et de
$-0{,}028$ sur les biais. C'est ce que fait aujourd'hui \texttt{generate.py}
(\texttt{src/stack.py}). Les poids sont figés dans le code, pour que chaque modèle ne soit
calculé qu'une fois sur la plateforme : 4,6 min sur entraînement + test, pour une limite de 20 min.}

\section{Que donnera la plateforme ?}

La plateforme donne à \texttt{generate.py} deux fois plus de notes (entraînement + test).
Pour l'estimer, on a recalculé la FM sur l'entraînement plus la moitié du test, et mesuré
la RMSE sur l'autre moitié (puis l'inverse) (\texttt{experiments/run\_platform\_proxy.py}).
Avec 50\,\% de notes en plus, la RMSE de la FM passe de 0,855 à 0,840. En prolongeant cette
tendance, on estime la FM seule autour de \textbf{0,83} sur la plateforme, et l'ensemble un
peu en dessous. Les scores de la plateforme et du test local ne sont pas directement
comparables : les notes cachées ne sont pas les mêmes.

%=============================================================================
\part{Récapitulatif}
%=============================================================================

\section{Toutes les méthodes}

\begin{center}
\begin{tabularx}{\linewidth}{@{}rXrr@{}}
\toprule
\textbf{Ordre} & \textbf{Méthode} & \textbf{RMSE test} & \textbf{Temps} \\
\midrule
1 & Moyenne globale & 1,0369 & 0 s \\
2 & Biais (référence) & 0,8746 & 0,01 s \\
3 & Factorisation : ALS / ALS-WR / SGD & 0,8725 / 0,8746 / 0,8775 & 0,5 à 1,5 s \\
4 & Soft-Impute & 0,8716 & 4,7 s \\
5 & Plus proches voisins (kNN) & 0,8705 & 1,1 s \\
6 & Biais + genres & 0,8736 & 0,04 s \\
7 & Mélange des méthodes 2 à 6 & 0,8683 & 7,7 s \\
8 & AutoRec & 0,8722 & 17 s \\
9 & BPMF & 0,8646 & 44 s \\
10 & FM bayésienne & 0,8552 & 60 s \\
11 & GC-MC, GLocal-K (réseaux de neurones) & 0,8792 / 0,8965 & 20 à 240 s \\
12 & \textbf{Ensemble final} (FM, FM corrigée, FM$_2$, $k$NN) & \textbf{0,8470} & 3 min \\
\bottomrule
\end{tabularx}
\end{center}

Tout a tourné sur un ordinateur portable, sans carte graphique. Chaque expérience n'a été
lancée qu'une fois, avec un seul tirage de la validation : des écarts de l'ordre de 0,001
entre deux méthodes ne sont pas forcément significatifs.

\section{Où est le code}

\begin{center}
\small
\begin{tabularx}{\linewidth}{@{}lX@{}}
\toprule
\textbf{Méthode} & \textbf{Code} et \textbf{script d'expérience} (résultats dans \texttt{results/}) \\
\midrule
Lecture des données, RMSE & \texttt{src/data.py} \\
Moyenne, biais & \texttt{src/baselines.py} \\
Factorisation & \texttt{src/mf.py} ; \texttt{experiments/run\_mf.py} \\
Soft-Impute, kNN, biais + genres, mélange & \texttt{src/soft\_impute.py}, \texttt{src/knn.py},
  \texttt{src/side\_info.py}, \texttt{src/blend.py} ; \texttt{experiments/run\_alternatives.py} \\
AutoRec & \texttt{src/autorec.py} ; \texttt{experiments/run\_autorec.py} \\
BPMF, FM bayésienne & \texttt{src/bpmf.py}, \texttt{src/bfm.py} ; \texttt{experiments/run\_bayesian.py} \\
Variantes de la FM & \texttt{src/bfm.py}, \texttt{src/bfm\_ordinal.py} ; \texttt{experiments/run\_fm\_variants.py} \\
GC-MC, GLocal-K & \texttt{src/gcmc.py}, \texttt{src/glocalk.py} ; \texttt{experiments/run\_literature\_models.py} \\
EASE, corrections des résidus & \texttt{src/ease.py}, \texttt{src/residual.py}, \texttt{src/residual\_gbdt.py} ;
  \texttt{experiments/run\_residual\_models.py} \\
Ensembles & \texttt{src/ensemble.py}, \texttt{src/stack.py} ; \texttt{experiments/run\_ensemble.py},
  \texttt{experiments/run\_stack.py} \\
Estimation du score plateforme & \texttt{experiments/run\_platform\_proxy.py} \\
Description des données & \texttt{experiments/eda.py} (figures dans \texttt{figures/eda/}) \\
\bottomrule
\end{tabularx}
\end{center}

\end{document}
